% ============================================================
%  CHAPTER 5 — Perception System
% ============================================================
\chapter{Perception System}
\label{ch:perception}

\section{Overview}

The perception system is responsible for providing the agent with an understanding
of the visual environment.
It has two modes of operation: a \textbf{detection mode} that runs YOLO on demand
when the agent is processing a command, and a \textbf{streaming mode} used by the
live visualiser to annotate the camera feed continuously.

The main component is \texttt{PerceptionManager} (v2), which abstracts:
\begin{itemize}[noitemsep]
  \item RGB image acquisition from the ROS camera topic.
  \item Object detection via the YOLO11n HTTP microservice.
  \item A local YOLOv4-tiny fallback if the service is unreachable.
  \item Detection visualisation (bounding boxes, labels, confidence scores).
\end{itemize}

\section{RGB-D Image Pipeline}

\subsection{Synchronised Acquisition}

RGB and depth images are acquired as time-synchronised pairs using
\texttt{message\_fil\-ters.\-Approx\-imate\-Time\-Syn\-chron\-izer} with a 100\,ms time tolerance:

\begin{lstlisting}[language=Python, caption={RGB-D synchronisation setup.}]
rgb_sub   = message_filters.Subscriber('/xtion/rgb/image_raw', Image)
depth_sub = message_filters.Subscriber('/xtion/depth_registered/image_raw', Image)
self.ts   = message_filters.ApproximateTimeSynchronizer(
    [rgb_sub, depth_sub], queue_size=5, slop=0.1)
self.ts.registerCallback(self._sync_callback)
\end{lstlisting}

This guarantees that the RGB and depth images used together for 3-D localisation
correspond to the same physical moment, eliminating stereo artefacts caused by
robot motion between frames.

\subsection{Depth Decoding}

The Xtion delivers depth in two possible encodings.
The callback handles both:

\begin{lstlisting}[language=Python, caption={Depth image decoding.}]
if depth_msg.encoding == '32FC1':
    arr = np.frombuffer(depth_msg.data, dtype=np.float32)
elif depth_msg.encoding == '16UC1':
    arr = np.frombuffer(depth_msg.data, dtype=np.uint16) \
          .astype(np.float32) / 1000.0  # mm -> m
self.latest_depth = arr.reshape(depth_msg.height, depth_msg.width)
\end{lstlisting}

Valid depth range is $[0.15, 3.5]$ m; values outside this range are set to NaN.

\section{YOLO Model Selection}

\subsection{Evolution of Models Used}

Four YOLO variants were evaluated during development.
Table~\ref{tab:yolo_eval} summarises the key metrics.

\begin{table}[H]
\centering
\caption{YOLO model evaluation on the deployment platform.}
\label{tab:yolo_eval}
\small
\begin{tabular}{llllll}
\toprule
\textbf{Model} & \textbf{Deployment} & \textbf{Infer.} & \textbf{Classes} & \textbf{Seg.} & \textbf{Decision} \\
\midrule
YOLOv4-tiny   & Docker (cv2.dnn)   & $\sim$50 ms  & 80 (fixed) & No  & Early v. \\
YOLO-World    & Host (Py 3.10)     & $\sim$150 ms & Open vocab & No  & Too slow \\
YOLOE-26s-seg & Host (Py 3.10)     & $\sim$300 ms & Open vocab & Yes & Too slow \\
YOLO11n       & Host microservice  & $\sim$18 ms  & 80 (COCO)  & No  & \textbf{Chosen} \\
\bottomrule
\end{tabular}
\normalsize
\end{table}

\subsection{YOLO11n}

YOLO11n is the nano variant of the Ultralytics YOLO11 architecture.
Key properties relevant to this deployment:

\begin{itemize}[noitemsep]
  \item Trained on MS-COCO 80 classes, covering all objects relevant to the scenario
        (\texttt{bottle}, \texttt{person}, \texttt{cup}, \texttt{chair}, \textit{etc.}).
  \item Single-stage detector: one forward pass for the full image.
  \item CPU inference at $\approx$18 ms per frame → theoretical 55 fps
        (limited in practice by the HTTP overhead to $\approx$15 fps).
  \item Weights size: $\approx$5.4 MB.
  \item Confidence threshold: 0.4 (empirically tuned to avoid false positives
        while not missing partially occluded objects).
\end{itemize}

\subsection{Why Not Open-Vocabulary Models?}

YOLO-World and YOLOE-26s-seg were attractive because they can detect \emph{any}
object described by a text prompt, rather than being limited to COCO classes.
However:

\begin{itemize}[noitemsep]
  \item \textbf{YOLO-World}: $\sim$150 ms CPU inference, acceptable for periodic
        detection but inadequate for the live 15 fps stream required by the visualiser.
  \item \textbf{YOLOE-26s-seg}: $\sim$300 ms CPU inference; additionally requires
        downloading the 242 MB mobileclip2\_b.ts CLIP model at startup, adding
        30--60 seconds of initialisation time.
\end{itemize}

For a future deployment with GPU access, open-vocabulary models would reduce inference
to $<$10 ms and would be the preferred choice.

\section{YOLO HTTP Microservice}

\subsection{Architecture}

Because the Docker container runs Python 3.6, which cannot use Ultralytics,
YOLO11n is deployed on the host machine as a standalone HTTP server
(\texttt{yolo\_service.py}):

\begin{figure}[H]
\centering
\begin{tikzpicture}[node distance=0.5cm and 2.5cm]
  \node[box, minimum width=3.5cm] (docker) {Docker Container\\PerceptionManager v2\\(Python 3.6)};
  \node[box, minimum width=3.5cm, right=3cm of docker] (host) {Host Machine\\YOLO11n Service\\(Python 3.10, port 5001)};

  \draw[arr, bend left=20] (docker.north east)
    to node[above, font=\small, align=center]{POST /detect\\JPEG bytes + X-Classes header} (host.north west);
  \draw[arr, bend left=20] (host.south west)
    to node[below, font=\small]{JSON detections} (docker.south east);
\end{tikzpicture}
\caption{YOLO HTTP microservice communication.}
\label{fig:yolo_service}
\end{figure}

\subsection{API}

\begin{table}[H]
\centering
\caption{YOLO service API endpoints.}
\label{tab:yolo_api}
\begin{tabularx}{\linewidth}{llX}
\toprule
\textbf{Method} & \textbf{Endpoint} & \textbf{Description} \\
\midrule
POST & \texttt{/detect} & Body: raw JPEG bytes. Optional header \texttt{X-Classes}: comma-separated class filter. Returns JSON detections. \\
GET  & \texttt{/health} & Returns model name, inference capability status. \\
\bottomrule
\end{tabularx}
\end{table}

The \texttt{X-Classes} header allows the caller to filter detections server-side,
reducing the JSON payload for targeted detection:

\begin{lstlisting}[language=Python, caption={Class-filtered detection request.}]
resp = requests.post(
    'http://localhost:5001/detect',
    data=jpeg_bytes,
    headers={
        'Content-Type': 'image/jpeg',
        'X-Classes': 'bottle'      # only return bottle detections
    },
    timeout=3.0
)
\end{lstlisting}

\subsection{Threaded Server}

The HTTP server uses \texttt{socketserver.ThreadingMixIn} so that MJPEG stream
connections and \texttt{/detect} requests are served concurrently without blocking:

\begin{lstlisting}[language=Python, caption={Threaded HTTP server.}]
class _ThreadedHTTPServer(socketserver.ThreadingMixIn, HTTPServer):
    daemon_threads = True
\end{lstlisting}

Without threading, a long-running \texttt{/detect} call would block the
\texttt{/health} endpoint, causing false negatives in service health checks.

\subsection{Detection Response Format}

\begin{lstlisting}[language=JSON, caption={YOLO service JSON response.}]
{
  "detections": [
    {
      "class_name": "bottle",
      "confidence": 0.8712,
      "bbox": [142, 201, 85, 178]
    }
  ],
  "model": "yolo11n",
  "inference_ms": 18,
  "segmentation": false
}
\end{lstlisting}

Bounding boxes are in $[x, y, w, h]$ pixel format where $(x,y)$ is the top-left corner.

\section{Live Visualiser}

\subsection{Web Dashboard (port 8080)}

\texttt{web\_visualiser.py} serves an MJPEG stream and JSON status endpoint at port 8080,
accessible from any browser on the network.
It subscribes to two ROS topics:

\begin{itemize}[noitemsep]
  \item \texttt{/xtion/rgb/image\_raw} — for continuous raw camera feed (always active).
  \item \texttt{/agent/camera\_annotated} — annotated frames published by the agent
        during active processing (overrides raw feed).
\end{itemize}

A key fix was adding \texttt{socketserver.ThreadingMixIn} to allow concurrent
connections: a single-threaded server caused the \texttt{/status} poll (every 500 ms)
to block behind the MJPEG stream, freezing the dashboard.

\subsection{Live YOLO Detector (port 8081)}

\texttt{live\_detector.py} provides a browser-accessible stream at port 8081 showing
real-time YOLO detections overlaid on the camera feed.
A \textbf{decoupled two-thread design} prevents slow detection from degrading stream
smoothness:

\begin{figure}[H]
\centering
\begin{tikzpicture}[node distance=0.4cm and 2.0cm]
  \node[box, minimum width=4cm] (cam)  {Camera Thread\\15 fps from /xtion};
  \node[box, minimum width=4cm, right=2cm of cam] (det) {Detection Thread\\calls YOLO service};
  \node[box, minimum width=4cm, below=1cm of cam] (stream){MJPEG Stream\\15 fps (browser)};
  \node[box, minimum width=4cm, below=1cm of det] (dets){Last known\\detections};

  \draw[arr] (cam.south) -- node[left, font=\small]{draws boxes} (stream.north);
  \draw[arr] (det.south) -- (dets.north);
  \draw[arr] (dets.west) -- node[below, font=\small]{shared state} (stream.east);
  \draw[arr, dashed] (cam.east) -- node[above, font=\small]{latest JPEG} (det.west);
\end{tikzpicture}
\caption{Decoupled architecture of the live YOLO detector.}
\label{fig:live_detector}
\end{figure}

The camera thread reads frames and overlays the \emph{most recently computed} detection
boxes, regardless of when the detection ran.
The detection thread calls the YOLO service as fast as possible (typically 8--15 Hz
for YOLO11n) and updates the shared detection list.
This ensures the stream stays at 15 fps even when detection takes 100+ ms.

\section{Local Fallback: YOLOv4-tiny}

If the YOLO service at \texttt{http://localhost:5001} is unreachable at startup,
\texttt{Perception\-Manager} falls back to a local YOLOv4-tiny model loaded via
\texttt{cv2.dnn.\-read\-Net\-From\-Darknet()}.
This requires the weights and configuration files to be present in the workspace
directory.
The fallback is slower and less accurate than YOLO11n but ensures the system can
still operate in degraded mode without network connectivity to the host.

\section{PerceptionManager Interface}

The public interface of \texttt{PerceptionManager} used by the agent and skills:

\begin{lstlisting}[language=Python, caption={PerceptionManager public API.}]
class PerceptionManager:
    def get_latest_rgb(self) -> np.ndarray:
        """Return latest BGR image from camera."""

    def run_yolo(self, image, target_classes=None) -> List[Dict]:
        """Run YOLO detection. Returns list of detections:
        [{'class_name': str, 'confidence': float, 'bbox': [x,y,w,h]}, ...]"""

    def detect_objects(self, target_classes=None) -> List[Dict]:
        """Capture a fresh image and run YOLO."""

    def draw_detections(self, image, detections) -> np.ndarray:
        """Draw bounding boxes and labels on image copy."""
\end{lstlisting}

The \texttt{target\_classes} parameter is forwarded as the \texttt{X-Classes} header
to the YOLO service, enabling server-side filtering without changing the detection model.
